% ===========================================================================
% Section 5 — Solution Approach / Methodology
% IEEE Conference / Journal Format
% ===========================================================================

\section{Solution Approach / Methodology}
\label{sec:methodology}

This paper proposes the \textbf{SHARP-LLM} (Secure Healthcare-Aligned Risk
Prediction via Lightweight LLMs) framework for automated source code
vulnerability detection and CWE classification.  The framework leverages
fine-tuned lightweight large language models operating on tokenised C/C++
source code to classify vulnerabilities into 118~CWE categories.  The
novelty of the approach is threefold: (i)~a \emph{template-aware data
splitting} strategy that eliminates structural data leakage inherent in the
Juliet Test Suite; (ii)~a \emph{multi-paradigm model evaluation} that
systematically compares encoder-based fine-tuning, QLoRA-based
parameter-efficient fine-tuning, and zero-shot generative inference within a
single unified pipeline; and (iii)~a healthcare-oriented risk
prioritisation layer that maps detected CWE categories to privacy, security,
and policy-relevant control domains.

% ---------------------------------------------------------------------------
\subsection{Pipeline Overview}
\label{sec:pipeline}

Fig.~\ref{fig:system_pipeline} presents the end-to-end SHARP-LLM pipeline.
Raw C/C++ source code first undergoes preprocessing and sub-word
tokenisation.  The tokenised representation is then fed into one of several
fine-tuned lightweight LLMs.  The model produces a binary safe/vulnerable
decision; for vulnerable inputs the system further performs multi-class CWE
association across 118~categories, maps the identified CWE to privacy and
security risk categories, aligns the risk with policy-relevant control
domains, and finally outputs a healthcare-oriented risk prioritisation
score.

\begin{figure}[t]
\centering
\resizebox{\columnwidth}{!}{%
\begin{tikzpicture}[
    node distance=4mm and 6mm,
    block/.style={
        draw,
        rounded corners,
        align=center,
        minimum width=2.6cm,
        minimum height=0.7cm,
        font=\footnotesize
    },
    line/.style={-Latex, thin}
]

\node[block] (input) {Source Code};
\node[block, below=of input] (prep) {Preprocessing and Tokenization};
\node[block, below=of prep] (llm) {Fine-Tuned Lightweight LLM};

\node[block, below left=6mm and 10mm of llm] (safe) {Safe Code};
\node[block, below right=6mm and 10mm of llm] (vuln) {Vulnerable Code};

\node[block, below=of vuln] (cwe) {CWE Association};
\node[block, below=of cwe] (risk) {Privacy and Security\\Risk Categorization};
\node[block, below=of risk] (policy) {Policy-Relevant\\Control Domains};
\node[block, below=of policy] (output) {Healthcare-Oriented\\Risk Prioritization Output};

\draw[line] (input) -- (prep);
\draw[line] (prep) -- (llm);
\draw[line] (llm) -- (safe);
\draw[line] (llm) -- (vuln);
\draw[line] (vuln) -- (cwe);
\draw[line] (cwe) -- (risk);
\draw[line] (risk) -- (policy);
\draw[line] (policy) -- (output);

\end{tikzpicture}%
}
\caption{High-level workflow of the proposed SHARP-LLM framework.}
\label{fig:system_pipeline}
\end{figure}

% ---------------------------------------------------------------------------
\subsection{Data and Environment Design}
\label{sec:data_design}

\subsubsection{Dataset}
The study employs the \textbf{NIST Juliet Test Suite for C/C++ v1.3}~\cite{juliet2017}, a widely used benchmark containing synthetic but
structurally realistic test cases spanning 118~Common Weakness Enumeration
(CWE) categories.  Each source file follows a naming convention encoding the
CWE identifier, functional variant, and flow variant, which we parse using a
regular-expression extractor adapted from the suite's \texttt{py\_common.py}
utility.  The resulting dataset comprises individual \texttt{.c}/\texttt{.cpp}
files, each labelled by CWE category and template identifier.

\subsubsection{Input Representation and Feature Space}
Each sample is a raw C/C++ source code string.  Preprocessing applies a
four-stage pipeline:
\begin{enumerate}
    \item \textbf{Header removal} --- strips the auto-generated
          \texttt{TEMPLATE GENERATED TESTCASE FILE} comment block.
    \item \textbf{Comment stripping} --- a character-level state machine
          removes single-line (\texttt{//}) and multi-line
          (\texttt{/* ... */}) comments while preserving string and
          character literals.
    \item \textbf{Guard removal} --- eliminates Juliet-specific preprocessor
          guards (\texttt{\#ifndef OMITBAD}, \texttt{\#ifndef OMITGOOD},
          \texttt{\#ifdef INCLUDEMAIN}, and trailing \texttt{\#endif}).
    \item \textbf{Whitespace normalisation} --- collapses consecutive
          whitespace and blank lines to produce a compact, semantics-preserving
          representation.
\end{enumerate}
The preprocessed code is tokenised on-the-fly using each model's native
sub-word tokeniser (SentencePiece for CodeT5 family, byte-pair encoding for
CodeBERT/GraphCodeBERT, Gemma tokeniser for CodeGemma), with a maximum
sequence length of $L = 256$ tokens, right-truncation, and
\texttt{max\_length} padding.

\subsubsection{Output Space}
The output is a 118-dimensional probability vector $\hat{\mathbf{y}} \in
\mathbb{R}^{118}$ produced via a softmax layer.  Each dimension corresponds
to one CWE category, mapped through a deterministic label encoding (CWE IDs
sorted numerically $\mapsto$ integer labels $0 \ldots 117$).  At inference
time the top-$k$ predictions with associated confidence scores are returned.

\subsubsection{Template-Aware Data Splitting}
A key design decision is the \emph{template-aware splitting} strategy.  In
the Juliet suite, each CWE category contains multiple \emph{templates}
(functional variants) instantiated across several \emph{flow variants}.
Flow variants of the same template share identical vulnerability logic,
differing only in control-flow structure.  A na\"ive random split risks
placing flow variants of the same template in both train and test sets,
causing structural data leakage and artificially inflated metrics.

We employ \texttt{GroupShuffleSplit} from scikit-learn with the template
identifier as the grouping key, performed independently within each CWE
category to maintain class balance:
\begin{equation}
    \forall\, t \in \mathcal{T}: \quad t \in \mathcal{D}_{\text{train}}
    \;\oplus\; t \in \mathcal{D}_{\text{test}}
\label{eq:template_split}
\end{equation}
where $\mathcal{T}$ is the set of all template identifiers and $\oplus$
denotes exclusive disjunction.  CWE categories with only a single template
are assigned entirely to the training set.  The split ratio is 80/20
(train/test) with a fixed random seed ($s = 42$) for reproducibility.

% ---------------------------------------------------------------------------
\subsection{Model Architectures}
\label{sec:model}

The framework evaluates three distinct model paradigms, enabling a
systematic comparison of fine-tuning strategies under resource-constrained
environments (8\,GB VRAM).

\subsubsection{Encoder-Based Classification}
Four pre-trained encoder models are fine-tuned with a classification head:

\begin{itemize}
    \item \textbf{CodeT5-Small} (60\,M params) and \textbf{CodeT5-Base}
          (220\,M params)~\cite{wang2021codet5}: T5-family encoder-decoder
          models pre-trained on CodeSearchNet with bimodal dual generation
          and identifier-aware objectives.  We use only the encoder.
    \item \textbf{CodeBERT-Base} (125\,M params)~\cite{feng2020codebert}:
          RoBERTa-based model pre-trained on programming and natural
          language with masked language modelling and replaced token
          detection.
    \item \textbf{GraphCodeBERT-Base} (125\,M params)~\cite{guo2020graphcodebert}:
          extends CodeBERT with data-flow graph pre-training objectives.
\end{itemize}

The classification architecture is:
\begin{equation}
    \hat{\mathbf{y}} = \text{softmax}\!\Big(
        \mathbf{W} \cdot \text{Dropout}\!\big(
            \bar{\mathbf{h}}
        \big) + \mathbf{b}
    \Big)
\label{eq:classifier}
\end{equation}
where $\bar{\mathbf{h}}$ is the mean-pooled encoder output over non-padding
tokens:
\begin{equation}
    \bar{\mathbf{h}} = \frac{
        \sum_{i=1}^{L} m_i \cdot \mathbf{h}_i
    }{
        \sum_{i=1}^{L} m_i
    }
\label{eq:mean_pool}
\end{equation}
with $\mathbf{h}_i \in \mathbb{R}^{d}$ being the hidden state at position
$i$, $m_i \in \{0,1\}$ the attention mask, and $d$ the model's hidden
dimension.

\subsubsection{QLoRA-Based Parameter-Efficient Fine-Tuning}
For larger decoder-only models, we employ Quantised Low-Rank Adaptation
(QLoRA)~\cite{dettmers2023qlora} applied to \textbf{CodeGemma-2B}
(v1.1, 2.5\,B parameters)~\cite{team2024codegemma}, a model
purpose-built for code tasks with dependency-graph and unit-test
lexical packing during pre-training.

The base model is loaded with 4-bit NormalFloat (NF4) quantisation via
\texttt{BitsAndBytesConfig} with double quantisation enabled, reducing
memory from $\approx\!5.5$\,GB (FP16) to $\approx\!1.8$\,GB.  Low-rank
adapters are applied to all linear projection layers:
\begin{equation}
    \mathbf{W}' = \mathbf{W}_{\text{frozen}}^{\text{NF4}}
                 + \frac{\alpha}{r}\,\mathbf{B}\mathbf{A}
\label{eq:lora}
\end{equation}
where $\mathbf{W}_{\text{frozen}}^{\text{NF4}}$ is the 4-bit quantised
weight matrix, $\mathbf{B} \in \mathbb{R}^{d \times r}$ and
$\mathbf{A} \in \mathbb{R}^{r \times k}$ are the low-rank adapter
matrices with rank $r = 16$ and scaling factor $\alpha = 32$.  Target
modules include query, key, value, output, gate, up, and down projections.
The classification head (\texttt{score} layer) is trained in full
precision via \texttt{modules\_to\_save}.

Table~\ref{tab:qlora_config} summarises the QLoRA configuration.

\begin{table}[t]
\centering
\caption{QLoRA configuration for CodeGemma-2B fine-tuning.}
\label{tab:qlora_config}
\begin{tabular}{ll}
\hline
\textbf{Parameter}        & \textbf{Value} \\
\hline
LoRA rank ($r$)            & 16 \\
LoRA alpha ($\alpha$)      & 32 \\
LoRA dropout               & 0.05 \\
Target modules             & q, k, v, o, gate, up, down \\
Quantisation type          & NF4 (double quant) \\
Compute dtype              & bfloat16 \\
Gradient checkpointing     & Enabled \\
Optimiser                  & PagedAdamW (8-bit) \\
Learning rate              & $2 \times 10^{-4}$ \\
Effective batch size       & 8 ($2 \times 4$ accum.) \\
Epochs                     & 5 (patience = 2) \\
\hline
\end{tabular}
\end{table}

\subsubsection{Zero-Shot Generative Inference}
Two generative models are evaluated without any training:
\textbf{CodeGemma-2B}~\cite{team2024codegemma} (completion mode) and
\textbf{Gemma-4-E2B-IT}~\cite{team2024gemma} (instruction-tuned mode, 5.1\,B
total parameters, 2.3\,B effective).

Zero-shot inference uses a structured prompt that enumerates all 118~valid
CWE identifiers and instructs the model to output only the CWE number.
For instruction-tuned models, prompts are wrapped with the model's chat
template.  Beam search ($B = 10$ for CodeGemma, $B = 2$ for Gemma~4~IT) is
used for decoding, with a regex-based parser extracting valid CWE numbers
from each beam.  Confidence scores are derived via softmax normalisation
of beam log-probabilities:
\begin{equation}
    c_j = \frac{\exp(s_j)}{\sum_{j'=1}^{K} \exp(s_{j'})}
\label{eq:beam_confidence}
\end{equation}
where $s_j$ is the log-probability of the $j$-th valid beam and $K$ is the
number of unique valid predictions.

% ---------------------------------------------------------------------------
\subsection{Novel Components}
\label{sec:novel}

\subsubsection{Template-Aware Anti-Leakage Splitting}
Existing studies on the Juliet Test Suite often perform random sample-level
splitting, which introduces template leakage --- flow variants of the same
template appearing in both train and test sets.  Since flow variants share
identical vulnerability logic, this inflates reported metrics.  Our
template-aware splitting (Section~\ref{sec:data_design}) enforces strict
separation at the template level, providing a more rigorous evaluation.

\subsubsection{Multi-Paradigm Comparative Framework}
Rather than evaluating a single model or fine-tuning strategy, SHARP-LLM
provides a unified pipeline for comparing three paradigms ---
encoder-based full fine-tuning, QLoRA parameter-efficient fine-tuning, and
zero-shot generative inference --- under identical data conditions,
preprocessing, and evaluation metrics.  This enables controlled ablation
of model architecture and training strategy effects.

\subsubsection{Resource-Constrained QLoRA Integration}
The QLoRA fine-tuning of CodeGemma-2B is specifically engineered for
consumer-grade hardware (8\,GB VRAM), demonstrating that decoder-based
models with 2.5\,B parameters can be fine-tuned for multi-class
vulnerability classification on a single GPU.  The combination of NF4
double quantisation, gradient checkpointing, 8-bit paged optimiser, and
adapter-only training reduces peak VRAM consumption to
$\approx\!3$--$4$\,GB.

% ---------------------------------------------------------------------------
\subsection{Training Procedure}
\label{sec:training}

\subsubsection{Encoder Model Training}
All encoder models are trained with the AdamW optimiser~\cite{loshchilov2019decoupled} using a learning
rate of $5 \times 10^{-5}$, weight decay of $0.01$, and dropout
$p = 0.1$.  The loss function is the standard cross-entropy:
\begin{equation}
    \mathcal{L} = -\frac{1}{N}\sum_{i=1}^{N}\sum_{c=1}^{C}
        y_{i,c} \log \hat{y}_{i,c}
\label{eq:ce_loss}
\end{equation}
where $N$ is the batch size and $C = 118$ is the number of classes.
Mixed-precision training (FP16) with gradient scaling is enabled on CUDA
devices.  Early stopping monitors the macro-F1 score on the test set with a
patience of 1~epoch.

\subsubsection{QLoRA Training}
The QLoRA training loop follows a similar structure but operates only on
the trainable parameters (LoRA adapters and the classification head,
$\approx\!0.5\%$ of total parameters).  The 8-bit PagedAdamW optimiser
is used with a higher learning rate ($2 \times 10^{-4}$) appropriate for
adapter tuning.  Mixed-precision uses bfloat16 (preferred for NF4 QLoRA
stability).  Gradient accumulation over 4~steps yields an effective batch
size of $8$.  Early stopping patience is set to 2~epochs, and the maximum
number of epochs is $5$.

\subsubsection{Evaluation Metrics}
All models are evaluated using the following macro-averaged metrics:
\begin{itemize}
    \item \textbf{Accuracy}: $\frac{1}{N}\sum_{i=1}^{N}
          \mathbb{1}[\hat{y}_i = y_i]$
    \item \textbf{Macro Precision, Recall, F1-Score}: averaged
          across all 118 classes with zero-division fallback
    \item \textbf{Matthews Correlation Coefficient (MCC)}: a balanced
          measure robust to class imbalance
    \item \textbf{Macro False Positive Rate (FPR)}: computed per-class
          from the confusion matrix:
          $\text{FPR}_c = \frac{\text{FP}_c}{\text{FP}_c + \text{TN}_c}$,
          then averaged
    \item \textbf{Inference Latency}: average over 20 forward passes
          (5-run warmup, CUDA synchronised)
\end{itemize}

% ---------------------------------------------------------------------------
\subsection{Algorithm / Pseudocode}
\label{sec:pseudocode}

Algorithm~\ref{alg:training} presents the unified training procedure for
both encoder and QLoRA models.  Algorithm~\ref{alg:zero_shot} describes
the zero-shot inference procedure with beam-search decoding and CWE
extraction.

\begin{algorithm}[t]
\caption{Encoder / QLoRA Training Procedure}
\label{alg:training}
\begin{algorithmic}[1]
\REQUIRE Dataset $\mathcal{D}$, model $\mathcal{M}$, epochs $E$,
         patience $P$, accumulation steps $A$
\ENSURE Trained model with best adapter/checkpoint
\STATE $\mathcal{D}_{\text{train}}, \mathcal{D}_{\text{test}}
       \leftarrow \textsc{TemplateAwareSplit}(\mathcal{D})$
\STATE Preprocess all code samples via pipeline
\STATE Initialise $\mathcal{M}$ (load pre-trained weights)
\IF{QLoRA mode}
    \STATE Quantise $\mathcal{M}$ to NF4 (4-bit)
    \STATE Apply LoRA adapters ($r\!=\!16, \alpha\!=\!32$)
    \STATE Freeze all base parameters
    \STATE Enable gradient checkpointing
\ENDIF
\STATE $\text{best\_f1} \leftarrow 0$; $\text{patience\_ctr} \leftarrow 0$
\FOR{epoch $= 1$ \TO $E$}
    \STATE $\mathcal{M}.\text{train}()$
    \FOR{step, batch $\in$ $\mathcal{D}_{\text{train}}$}
        \STATE $\hat{\mathbf{y}} \leftarrow \mathcal{M}(\text{batch})$
        \STATE $\ell \leftarrow \textsc{CrossEntropy}(\hat{\mathbf{y}},
               \mathbf{y}) \,/\, A$
        \STATE $\ell.\text{backward}()$
        \IF{step $\bmod$ $A = 0$}
            \STATE Optimiser step; zero gradients
        \ENDIF
    \ENDFOR
    \STATE $\text{F1} \leftarrow \textsc{Evaluate}(\mathcal{M},
           \mathcal{D}_{\text{test}})$
    \IF{$\text{F1} > \text{best\_f1}$}
        \STATE $\text{best\_f1} \leftarrow \text{F1}$;
               save checkpoint
        \STATE $\text{patience\_ctr} \leftarrow 0$
    \ELSE
        \STATE $\text{patience\_ctr} \leftarrow
               \text{patience\_ctr} + 1$
    \ENDIF
    \IF{$\text{patience\_ctr} \geq P$}
        \STATE \textbf{break} \COMMENT{Early stopping}
    \ENDIF
\ENDFOR
\RETURN best checkpoint / adapter
\end{algorithmic}
\end{algorithm}

\begin{algorithm}[t]
\caption{Zero-Shot CWE Inference with Beam Search}
\label{alg:zero_shot}
\begin{algorithmic}[1]
\REQUIRE Code sample $x$, model $\mathcal{M}$, valid CWEs
         $\mathcal{V}$, beams $B$, top-$k$
\ENSURE Top-$k$ CWE predictions with confidence
\STATE $x' \leftarrow \textsc{Preprocess}(x)$
\STATE Truncate $x'$ to 800 characters
\STATE $p \leftarrow \textsc{BuildPrompt}(x', \mathcal{V})$
\IF{$\mathcal{M}$ is instruction-tuned}
    \STATE $p \leftarrow \textsc{ApplyChatTemplate}(p)$
\ENDIF
\STATE $\mathcal{S} \leftarrow \mathcal{M}.\text{generate}(p,
       \text{beams}\!=\!B)$
\STATE $\mathcal{C} \leftarrow \emptyset$
\FOR{sequence $s_j$, score $\sigma_j$ $\in$ $\mathcal{S}$}
    \STATE $\text{cwe} \leftarrow \textsc{RegexExtract}(s_j)$
    \IF{$\text{cwe} \in \mathcal{V}$ \AND $\text{cwe} \notin \mathcal{C}$}
        \STATE $\mathcal{C} \leftarrow \mathcal{C} \cup
               \{(\text{cwe}, \sigma_j)\}$
    \ENDIF
\ENDFOR
\STATE $c_j \leftarrow \textsc{Softmax}(\{\sigma_j\})$ for all
       $j \in \mathcal{C}$
\RETURN top-$k$ from $\mathcal{C}$ sorted by $c_j$
\end{algorithmic}
\end{algorithm}
